\documentclass[12pt]{article}
\usepackage[utf8]{inputenc}
\usepackage[T1]{fontenc}
\usepackage{hyperref}
\title{Sui Language Resources for K-12 Curriculum Development}
\author{Compiled by Claude Sonnet 5.5 (claude-sonnet-5-5) for Andr\'as Kornai}
\date{}
\begin{document}
\maketitle

\section{Introduction}
Sui (Shui, ISO 639-3 swi) is a Kam--Sui (Kra-Dai) language of Guizhou, China, with about 300,000 speakers. Most live in Sandu Shui Autonomous County, and some live in neighbouring Guangxi. The searches found general reference pages and information on the Shuishu script, but no school materials.

\section{Primary Educational Resources}
No Sui primers, textbooks or school programmes were found in the search results.

\section{Linguistic and Reference Materials}
\begin{itemize}
\item Wikipedia, Sui language: overview of speakers, dialects and writing systems (Latin, Sui script, Chinese characters). \url{https://en.wikipedia.org/wiki/Sui_language}
\item Wikipedia, Sui script (Shuishu): a logographic script used by shamans for divination and geomancy. \url{https://en.wikipedia.org/wiki/Sui_script}
\item Omniglot, Sui script. \url{https://www.omniglot.com/writing/sui.htm}
\item Letter from the Shui community to the Unicode working group (WG2 N4981): mentions the published Common Dictionary of Shuishu (水书常用字字典) and efforts to encode the script. \url{https://unicode.org/wg2/docs/n4981-LetterFromShuiCommunity.pdf}
\end{itemize}

\section{Religious and Community Resources}
\begin{itemize}
\item Asia Harvest people-group profile of the Shui (PDF). \url{https://www.people-groups.asiaharvest.org/China/chinaPeoples/S/Shui.pdf}
\end{itemize}

\section{Digital and Audio Resources}
No Sui audio, video or corpus resources were found. The Unicode letter above relates to digital encoding of the Shuishu script and mentions a 2006 input method.

\section{Curriculum Development Context}
Sui is not known to be taught in schools. The Shuishu script is not phonetic and was traditionally known to shamans. It was listed as Chinese intangible cultural heritage in 2006. The practical literacy options are Latin script or Chinese characters, and a curriculum would have to build most materials from scratch. A detailed review of the Shuishu dictionary was not found.

\end{document}